% ============================================================================
%  Capitolo 22 -- Verm1, parte 4: il gioco finito                (Lezione 22)
% ============================================================================
\chapter{Verm1, parte 4: il gioco finito}
\label{ch:verm1p4}

\epigraphline{L'ottimizzazione prematura è la radice di tutti i mali.}{Donald Knuth, 1974}

\begin{chapgoals}
Alla fine di questo capitolo avrete
\begin{itemize}
\item il punteggio e il punteggio migliore scritti in cifre, convertiti dalla vostra divisione per dieci;
\item un verme che va più veloce al crescere del punteggio, fino a un limite;
\item due suoni: una nota breve per un pasto, una lunga e grave per la fine;
\item una funzione \code{draw} che tocca una manciata di celle per passo invece dell'intero schermo;
\item \vermone finito, pronto per essere impacchettato nel Capitolo~\ref{ch:store}.
\end{itemize}
\end{chapgoals}

In quest'ultimo tutorial sul programma entra nel gioco la macchina stessa. La
frase di Knuth ne è il motto: il gioco è stato fatto \emph{funzionare} nella
parte~3, e solo ora, con una misura in mano, viene reso \emph{veloce}. La
Tabella~\ref{tab:v4:from} elenca ciò che si prende dai Capitoli da
\ref{ch:processor} a~\ref{ch:io}; torna anche il
Capitolo~\ref{ch:information}, perché una cifra diventa un carattere sommando
48, il codice ASCII di \texttt{0}. Il file di riferimento è
\fname{verm1\_part4.S} (Appendice~\ref{app:verm1}).

\begin{table}[htb]
\centering\small
\caption{Che cosa la parte 4 prende dai capitoli sulla macchina.}\label{tab:v4:from}
\begin{tabularx}{\textwidth}{@{}lLL@{}}
\toprule
\textbf{Da} & \textbf{Idea} & \textbf{Dove compare}\\
\midrule
Capitolo \ref{ch:processor} & la convenzione di chiamata: registri \reg{a}, \reg{t}, \reg{s}, \reg{ra}, lo stack & piccole funzioni con contratti; funzioni foglia\\
Capitolo \ref{ch:processor} & l'estensione M: \ins{divu} e \ins{remu}, quoziente e resto & le cifre del punteggio\\
Capitolo \ref{ch:instructions} & il ciclo del processore: ogni istruzione costa tempo & la misura prima e dopo il passo 4\\
Capitolo \ref{ch:video} & lo schermo è memoria: scrivere pixel è il costo & disegnare solo ciò che è cambiato\\
Capitolo \ref{ch:io} & i dispositivi appartengono al runtime: il programma chiede, il runtime fa & \code{prg32\_audio\_note}, con una durata\\
\bottomrule
\end{tabularx}
\end{table}

\begin{figure}[htb]
\centering
\begin{subfigure}[b]{0.32\textwidth}\includegraphics[width=\textwidth]{verm1_p4_title.png}\caption{\texttt{TITLE}, con il punteggio migliore}\end{subfigure}\hfill
\begin{subfigure}[b]{0.32\textwidth}\includegraphics[width=\textwidth]{verm1_p4_play.png}\caption{\texttt{PLAY}: il punteggio in cifre}\label{fig:v4:play}\end{subfigure}\hfill
\begin{subfigure}[b]{0.32\textwidth}\includegraphics[width=\textwidth]{verm1_p4_over.png}\caption{\texttt{OVER}}\label{fig:v4:over}\end{subfigure}
\caption{Dove va questo capitolo.}
\end{figure}

Partite da una parte~3 che funziona; se la vostra non funziona, copiate
\fname{verm1\_part3.S}. Il capitolo ha quattro passi, e il programma funziona
dopo ciascuno.

\begin{prgtry}[Prima di cominciare: misurate]
Eseguite la parte~3 in modalità debug e annotate ciò che la telemetria mostra
durante una partita: istruzioni completate, frame al secondo, frame in
ritardo. I numeri serviranno nel passo~5.
\end{prgtry}

% ----------------------------------------------------------------------------
\section{Passo 1: il punteggio in cifre}

\subsection{Un numero non è le sue cifre}

La variabile \code{score} contiene un numero binario: 137 è \hex{89}, un byte.
Lo schermo vuole tre caratteri, \texttt{'1'}, \texttt{'3'}, \texttt{'7'}, cioè
i byte \hex{31}, \hex{33}, \hex{37}. La conversione è l'algoritmo del
Capitolo~\ref{ch:integers}, eseguito dal programma invece che a mano: dividere
per dieci, tre volte.
\begin{align*}
137 / 10 &= 13, &&\text{resto 7: le unità;}\\
13 / 10 &= 1, &&\text{resto 3: le decine;}\\
1 / 10 &= 0, &&\text{resto 1: le centinaia.}
\end{align*}
Le cifre escono da destra, e sono memorizzate dall'ultimo carattere al primo.
Una cifra $d$ diventa un carattere sommando 48.

\subsection{\texttt{to\_decimal}}

\asmfile[caption={Passo 1 \textbf{(nuovo)}: \code{to\_decimal}.},label=lst:v4:dec]{code/l22/s08.S}

Il contratto del Listato~\ref{lst:v4:dec} è: $\reg{a0}=\text{valore}$,
$\reg{a1}=\text{dove vanno i tre caratteri}$. \ins{remu} dà la cifra e
\ins{divu} ciò che resta: due istruzioni dell'estensione M. \ins{sb} memorizza
un byte agli scostamenti 2, 1 e 0 da \reg{a1}. La funzione è una foglia che
usa solo registri \reg{t}: non le serve alcun frame. Valori sopra 999
mostrerebbero solo le ultime tre cifre, un limite che nessun giocatore
incontrerà presto. Le istruzioni di divisione sono lente, decine di cicli di
clock ciascuna su un processore piccolo; dieci per passo del verme non sono
niente.

\subsection{\texttt{draw\_score}: due numeri, due testi}

\asmfile[caption={Passo 1 \textbf{(modificato)}: \code{draw\_score}.},label=lst:v4:score]{code/l22/s09.S}

Sostituite la \code{draw\_score} della parte~3 con il
Listato~\ref{lst:v4:score}: le tacche spariscono. Le cifre vengono scritte in
una stringa in RAM, e poi la stringa viene disegnata come qualsiasi altra;
l'indirizzo \code{best\_txt+5} salta la parola \texttt{BEST} e lo spazio,
cinque caratteri. Il testo è disegnato con il suo colore di sfondo, così che
le nuove cifre coprono le vecchie senza alcuna cancellazione.

\subsection{Le variabili della parte 4}

\asmfile[caption={Passo 1 \textbf{(modificato)}: \code{.data}.},label=lst:v4:data]{code/l22/s10.S}

Sostituite la \code{.data} della parte~3 con il Listato~\ref{lst:v4:data}, e
scrivetela tutta ora. Le due stringhe \code{score\_txt} e \code{best\_txt}
sono in \code{.data} e non in \code{.rodata}, perché il programma vi scrive.
Le sei variabili da \code{dirty} a \code{erase\_y} non sono usate fino al
passo~3.

\begin{prgcheck}[Eseguite: le cifre]
\begin{enumerate}[nosep]
\item Compilate ed eseguite, premete \key{A}: la riga in alto dice
  \texttt{VERM1 000} e, a destra, \texttt{BEST 000}.
\item Giocate e mangiate: \texttt{001}, \texttt{002}, \ldots
\item Perdete: nella schermata di fine partita \texttt{BEST} mostra il vostro
  punteggio.
\item Nel debugger trovate \code{score\_txt} in memoria e leggetene i byte:
  \texttt{30 30 33 00} dopo tre pasti.
\end{enumerate}
\end{prgcheck}

% ----------------------------------------------------------------------------
\section{Passo 2: i livelli}

La velocità diventa una funzione del punteggio.

\asmfile[caption={Passo 2 \textbf{(modificato)}: le costanti. Solo le due righe \code{SPEED} differiscono dalla parte~3.},label=lst:v4:const]{code/l22/s13.S}

\code{SPEED\_Q8} è sostituita da un minimo e un massimo
(Listato~\ref{lst:v4:const}): $32/256$, un passo ogni 8 frame, all'inizio, e
$96/256$, tre passi ogni 8 frame, al massimo.

\asmfile[caption={Passo 2 \textbf{(modificato)}: il timer alla fine di \code{play\_update}.},label=lst:v4:speed]{code/l22/s14.S}

La velocità è $32+4\times\mathit{score}$: uno scorrimento e una somma
(Listato~\ref{lst:v4:speed}). Poi \ins{ble} e \ins{mv} prendono il minore tra
la velocità e 96. È una \emph{saturazione}: il valore segue una retta e poi
resta al limite, come la limitazione del Capitolo~\ref{ch:integers}. Il limite
è raggiunto a 16 punti (Tabella~\ref{tab:v4:levels}). Con la virgola fissa la
velocità cambia di un sessantaquattresimo di cella per frame a ogni pasto:
nessuna velocità intera potrebbe essere così fine.

\begin{table}[htb]
\centering\small
\caption{I livelli, in numeri. $256/\mathit{velocità}$ dà i frame per
passo.}\label{tab:v4:levels}
\begin{tabular}{@{}cccc@{}}
\toprule
\textbf{Punteggio} & \textbf{Velocità ($1/256$ cella/frame)} & \textbf{Frame per passo} & \textbf{Celle/s a 30 fps}\\
\midrule
0 & 32 & 8 & 3,75\\
4 & 48 & 5,3 & 5,6\\
8 & 64 & 4 & 7,5\\
12 & 80 & 3,2 & 9,4\\
16 e oltre & 96 & 2,7 & 11,25\\
\bottomrule
\end{tabular}
\end{table}

Le frazioni di frame sono reali: a $48/256$ i passi arrivano ogni 5 o 6 frame,
5,33 in media. Calcolate la riga del punteggio 2 per esercizio.

\begin{prgcheck}[Eseguite: sempre più veloce]
\begin{enumerate}[nosep]
\item Compilate ed eseguite: il verme parte più lentamente che nella parte~3,
  e acquista velocità a ogni pasto.
\item Contate con il debugger i frame tra due passi a punteggio 0: otto.
\item Arrivate a 16 punti: la velocità non cambia più.
\item Provate $\code{SPEED\_MAX}=256$, se ne avete il coraggio.
\end{enumerate}
\end{prgcheck}

% ----------------------------------------------------------------------------
\section{Passo 3: \texttt{step} dice che cosa ha cambiato}

Questo passo aggiunge i suoni, e le annotazioni che serviranno a \code{draw}
nel passo~4.

\subsection{\texttt{play\_note}: una funzione, qualsiasi nota}

\asmfile[caption={Passo 3 \textbf{(modificato)}: \code{play\_note} sostituisce \code{eat\_sound}.},label=lst:v4:note]{code/l22/s18.S}

Il chiamante sceglie la nota e la durata, in \reg{a2} e \reg{a4}; la funzione
riempie gli altri tre argomenti (Listato~\ref{lst:v4:note}). Gli argomenti
restano nei registri in cui \code{prg32\_audio\_note} li aspetta: progettare
un contratto attorno al chiamato fa risparmiare spostamenti.

\subsection{\texttt{step}: che cosa lascia a \texttt{draw}}

La funzione \code{step} della parte~3 viene modificata in cinque punti.

\asmfile[caption={Passo 3 \textbf{(modificato)}: \code{step} ricorda la vecchia testa.},label=lst:v4:s1]{code/l22/s19.S}

Si aggiungono due store (Listato~\ref{lst:v4:s1}): \code{old\_x} e
\code{old\_y}, la cella che la testa sta lasciando, che \code{draw} ridipingerà
nel colore del corpo.

\asmfile[caption={Passo 3 \textbf{(modificato)}: \code{step}, quando il verme mangia.},label=lst:v4:s2]{code/l22/s20.S}

Quando il verme mangia, la coda resta e non c'è niente da cancellare:
$\code{erase\_on}=0$ lo dice a \code{draw} (Listato~\ref{lst:v4:s2}).

\asmfile[caption={Passo 3 \textbf{(modificato)}: \code{step} annota la coda, e chiama la fine.},label=lst:v4:s3]{code/l22/s21.S}

Quando non mangia, la coda lascia una cella, e \code{step} ricorda quale:
$\code{erase\_on}=1$ significa che \code{draw} deve dipingerla di nero
(Listato~\ref{lst:v4:s3}). La fine di questo frammento è più corta che nella
parte~3, perché la collisione è ora una sola chiamata, a
\code{game\_is\_over}.

\asmfile[caption={Passo 3 \textbf{(modificato)}: \code{step} scrive la nuova testa (come nella parte~3).},label=lst:v4:s4]{code/l22/s22.S}

\asmfile[caption={Passo 3 \textbf{(nuovo)}: la fine di \code{step}.},label=lst:v4:s5]{code/l22/s23.S}

Alla fine (Listato~\ref{lst:v4:s5}) $\code{dirty}=1$ registra che è avvenuto
un passo e che \code{draw} ha del lavoro da fare. Un pasto è un punto, nuovo
cibo e una nota breve e acuta: la nota 72 per 60 millisecondi, il suono della
parte~2 attraverso la nuova funzione.

\subsection{\texttt{game\_is\_over}: la transizione, in un solo posto}

\asmfile[caption={Passo 3 \textbf{(nuovo)}: \code{game\_is\_over}.},label=lst:v4:over]{code/l22/s24.S}

Una transizione della macchina a stati merita una funzione tutta sua
(Listato~\ref{lst:v4:over}): lo stato, le uscite e il suono stanno in un solo
posto. La funzione memorizza \texttt{OVER}, aggiorna il punteggio migliore
come nella parte~3, pone $\code{dirty}=2$ perché va disegnata una nuova
schermata, e suona una nota lunga e grave: la 36, per 400 millisecondi.

\begin{prgcheck}[Eseguite: il suono della sconfitta]
\begin{enumerate}[nosep]
\item Compilate ed eseguite: il gioco ha esattamente l'aspetto che aveva dopo
  il passo~2.
\item Un pasto suona come prima; mordersi dà una nota lunga e grave.
\item Nel debugger guardate \code{erase\_x} ed \code{erase\_y} seguire la
  coda, un passo indietro.
\item \code{dirty} resta a 1 o 2: nessuno lo legge ancora.
\end{enumerate}
\textbf{Risultato atteso.} Nessun cambiamento visibile: lo schermo è ancora
disegnato per intero a ogni frame. Le nuove variabili sono scritte e non
ancora usate.
\end{prgcheck}

% ----------------------------------------------------------------------------
\section{Passo 4: disegnare solo ciò che è cambiato}

\subsection{Quanto costa un frame?}

La parte~3 cancella 64\,000 pixel e ridisegna il titolo, il punteggio,
l'intero verme e il cibo, trenta volte al secondo. Ma tra due passi sullo
schermo non cambia niente, e un passo cambia pochissimo: al più la cella che
la coda ha lasciato, che diventa nera; la vecchia testa, che diventa corpo; la
nuova testa; e forse il cibo e il punteggio.

\begin{prgconcept}
Ricordate che cosa è cambiato, e disegnate solo quello. L'immagine resta sullo
schermo da sola: il frame buffer è memoria (Capitolo~\ref{ch:video}).
\end{prgconcept}

Un verme di 200 celle costa 200 rettangoli per frame nella parte~3, e tre per
passo nella parte~4. È ciò che facevano i vecchi home computer, perché non
avevano scelta.

\subsection{\texttt{draw}: niente, poco, o tutto}

\asmfile[caption={Passo 4 \textbf{(modificato)}: \code{verm1\_draw}.},label=lst:v4:draw]{code/l22/s28.S}

La variabile \code{dirty} ha tre valori (Listato~\ref{lst:v4:draw}). Con 0 la
funzione ritorna subito, senza frame e senza chiamate: il caso comune, in cui
non è cambiato niente, costa cinque istruzioni macchina. Con 1 è avvenuto un
passo, e si disegnano i cambiamenti. Con 2 c'è una nuova schermata, e si
disegna tutto. L'indicatore viene azzerato appena è stato letto.

\subsection{\texttt{draw\_changes}: tre celle e il punteggio}

\asmfile[caption={Passo 4 \textbf{(nuovo)}: \code{draw\_changes}.},label=lst:v4:changes]{code/l22/s29.S}

L'ordine conta (Listato~\ref{lst:v4:changes}): prima la cella che la coda ha
lasciato, se c'è, in nero; poi la vecchia testa, nel colore del corpo; poi la
nuova testa, il cibo e il punteggio. Poiché la testa può entrare nella cella
che la coda ha appena lasciato, cancellare prima e disegnare la testa dopo dà
l'immagine giusta.

\subsection{\texttt{draw\_full}: un'intera schermata}

\asmfile[caption={Passo 4 \textbf{(nuovo)}: \code{draw\_full}, prima metà.},label=lst:v4:full1]{code/l22/s30.S}

\asmfile[caption={Passo 4 \textbf{(nuovo)}: \code{draw\_full}, \texttt{PLAY} e \texttt{OVER}.},label=lst:v4:full2]{code/l22/s31.S}

\code{draw\_full} è il corpo della vecchia \code{verm1\_draw}, ed è chiamata
una volta per ogni cambio di stato (Listati \ref{lst:v4:full1}
e~\ref{lst:v4:full2}). Il punteggio è disegnato anche nella schermata del
titolo.

\subsection{Tre piccole funzioni}

\asmfile[caption={Passo 4 \textbf{(nuovo)}: \code{draw\_head}.},label=lst:v4:head]{code/l22/s32.S}

\asmfile[caption={Passo 4 \textbf{(nuovo)}: \code{draw\_food}.},label=lst:v4:food]{code/l22/s33.S}

\asmfile[caption={Passo 4: \code{draw\_text}, la funzione di servizio della parte~3.},label=lst:v4:text]{code/l22/s34.S}

\subsection{Ogni transizione chiede una schermata}

Tre piccole aggiunte completano il meccanismo: dovunque cambi lo stato,
\code{dirty} diventa 2.

\asmfile[caption={Passo 4 \textbf{(modificato)}: \code{verm1\_init} chiede la prima schermata.},label=lst:v4:init]{code/l22/s35.S}

\asmfile[caption={Passo 4 \textbf{(modificato)}: \code{verm1\_update}, da \texttt{OVER} a \texttt{TITLE}.},label=lst:v4:upd]{code/l22/s36.S}

\asmfile[caption={Passo 4 \textbf{(modificato)}: la fine di \code{new\_game}, da \texttt{TITLE} a \texttt{PLAY}.},label=lst:v4:new]{code/l22/s37.S}

\begin{prgwarn}
Ogni transizione della macchina chiede un disegno completo; ogni passo ne
chiede uno parziale. Dimenticarne una è l'errore tipico di questa tecnica: lo
schermo mostra uno stato che non è più vero.
\end{prgwarn}

\begin{prgcheck}[Eseguite: Verm1 finito]
Lo stesso gioco che c'era dopo il passo~3, pixel per pixel (Figure
\ref{fig:v4:play} e~\ref{fig:v4:over}). Se vedete scie o buchi, un cambiamento
non è stato registrato: usate il debugger.
\end{prgcheck}

% ----------------------------------------------------------------------------
\section{Passo 5: misurare}

\prgqt è ora usato come uno strumento: contare prima di credere.

\begin{prgtry}[Prima e dopo, con i numeri]
\begin{enumerate}[nosep]
\item Eseguite il gioco finito in modalità debug e avviate una partita.
\item Leggete la telemetria durante dieci secondi di gioco: istruzioni
  completate, frame al secondo, frame in ritardo.
\item Confrontate con i numeri annotati per la parte~3 all'inizio del
  capitolo.
\item Lasciate crescere il verme fino a venti celle in entrambe le versioni e
  leggete di nuovo: un numero cresce con il verme, l'altro no.
\item Mettete un breakpoint sulla prima istruzione di \code{draw\_changes}: si
  ferma una volta per passo, non una volta per frame.
\item Scrivete i numeri nei vostri appunti: vanno nella descrizione del gioco
  nel Capitolo~\ref{ch:store}.
\end{enumerate}
\end{prgtry}

I numeri deliberatamente non sono stampati qui: vanno misurati. Un'avvertenza
nel leggerli: la telemetria conta le istruzioni della cartuccia, non quelle
del runtime che riempie i rettangoli, così che il risparmio reale sulla scheda
è maggiore di quanto il conteggio mostri; lì ogni pixel viaggia anche su SPI
verso il display.

\subsection{Seguire una chiamata con il debugger}

La convenzione di chiamata del Capitolo~\ref{ch:processor} si può guardare
istruzione per istruzione.

\begin{enumerate}
\item Fermatevi su \code{call to\_decimal} dentro \code{draw\_score}, e
  leggete \reg{a0}, \reg{a1}, \reg{ra} e \reg{sp}.
\item Fate un passo: il PC salta dentro \code{to\_decimal}, e \reg{ra}
  contiene ora l'indirizzo dell'istruzione dopo la chiamata.
\item Proseguite passo-passo: cambiano solo i registri \reg{t} e \reg{a0}, e
  \reg{sp} non si muove mai. È una funzione foglia.
\item Alla \ins{ret} il PC prende il valore di \reg{ra}: il programma è
  tornato in \code{draw\_score}.
\item Fate lo stesso con \code{call draw\_text}: \reg{sp} scende di 16,
  \reg{ra} è memorizzato a \code{12(sp)}, e ritorna.
\end{enumerate}

Nel listato del debugger \ins{call} compare come \ins{jal}, o come \ins{auipc}
seguita da \ins{jalr}, e \ins{ret} come \code{c.jr ra}.

% ----------------------------------------------------------------------------
\section{Verificate di aver capito}

\begin{quiz}{\code{to\_decimal} riceve $\reg{a0}=205$. Quali byte memorizza, dallo scostamento 0 al 2?}
\item \texttt{02 00 05}
\item \texttt{32 30 35}
\item \texttt{35 30 32}
\item \texttt{cd 00 00}
\end{quiz}

\begin{quiz}{A quale punteggio il verme fa un passo ogni 4 frame?}
\item 4
\item 8
\item 16
\item 32
\end{quiz}

\begin{quiz}{\code{dirty} vale 0 quando viene chiamata \code{verm1\_draw}. Quante istruzioni macchina esegue?}
\item 1
\item 5
\item circa 50
\item migliaia
\end{quiz}

\begin{quiz}{Perché \code{to\_decimal} può fare a meno di uno stack frame?}
\item è corta
\item non chiama funzioni e usa solo registri \reg{a} e \reg{t}
\item non ha argomenti
\item i frame sono facoltativi
\end{quiz}

\section{Fatelo vostro}

\begin{exercise}[su \prg]
\begin{enumerate}[nosep]
\item Una pausa: il pulsante \key{B} ferma il gioco e scrive \texttt{PAUSE};
  \key{B} di nuovo lo riprende. Un quarto stato.
\item Una melodia nella schermata del titolo: una tabella di note in
  \code{.rodata}, una nota ogni 15 frame.
\item Cibo bonus: un pezzo ogni cinque è giallo, vale tre punti, e scompare
  dopo 100 frame.
\item I vostri colori con \code{prg32\_palette\_set}, o le vostre regole:
  questo ora è il vostro gioco.
\end{enumerate}
\expected{La versione pubblicata nel Capitolo~\ref{ch:store} dovrebbe avere
qualcosa di vostro. Cambiate anche la stringa del titolo. (Suggerimenti: la
pausa è uno stato \texttt{PAUSED} con transizioni sul fronte di \key{B},
\hex{20}, e $\code{dirty}=2$ a entrambe le transizioni. Alla melodia servono
un indice e un contatore di frame in \code{.data}; la tabella si legge come
\code{dx}.)}
\end{exercise}

\section{Le funzioni che non sono cambiate}

Otto parti del file sono quelle della parte~3 e restano come sono. Sono
ristampate qui perché questo capitolo contenga l'intero programma.

\asmfile[caption={\textbf{Invariato}: \code{verm1\_update}, pulsanti e fronti.},label=lst:v4:u1]{code/l22/s48.S}
\asmfile[caption={\textbf{Invariato}: \code{new\_game}, la mappa e il verme.},label=lst:v4:u2]{code/l22/s49.S}
\asmfile[caption={\textbf{Invariato}: \code{play\_update}, la guida.},label=lst:v4:u3]{code/l22/s50.S}
\asmfile[caption={\textbf{Invariato}: \code{step}, la ricerca nella tabella.},label=lst:v4:u4]{code/l22/s51.S}
\asmfile[caption={\textbf{Invariato}: \code{place\_food}.},label=lst:v4:u5]{code/l22/s52.S}
\asmfile[caption={\textbf{Invariato}: \code{draw\_worm}, usata solo da \code{draw\_full}.},label=lst:v4:u6]{code/l22/s53.S}
\asmfile[caption={\textbf{Invariato}: \code{draw\_cell}.},label=lst:v4:u7]{code/l22/s54.S}
\asmfile[caption={\textbf{Invariato}: \code{.rodata}.},label=lst:v4:u8]{code/l22/s55.S}

\section{Dove siamo, e che cosa viene dopo}

\vermone è finito: circa seicento righe di assembly \riscv, tutte scritte da
voi. È cominciato come tre funzioni che ritornano, ed è cresciuto un passo
eseguibile alla volta. Dentro c'è ogni capitolo del libro: numeri, logica,
stato, memoria, chiamate, pixel, dispositivi.

Un programma che solo il suo autore può eseguire è finito a metà. Il prossimo
capitolo gli dà un nome, un'icona, i crediti, un bundle per due architetture,
e il Cartridge Store.
